\section{Implement a Stack using a Queue}
\label{sec:sq-stack-using-queue}

\problemstatement{Implement a stack (LIFO: \textit{push}, \textit{pop},
\textit{top}) using only the standard operations of a queue
(\textit{enqueue}, \textit{dequeue}, \textit{front}) as building blocks --
no direct array or linked-list indexing allowed.}

\begin{patternbox}
\emph{``Implement structure $X$ using only structure $Y$'s operations''}
is a request to simulate one access order using a tool built for the
opposite access order. The trick, here and in
Section~\ref{sec:sq-queue-using-stack}, is realizing that you can pay an
\textbf{extra cost at one specific operation} (either push or pop, your
choice) to physically reorder the underlying structure's contents,
restoring the access pattern you actually need.
\end{patternbox}

\subsection*{Understanding the problem carefully}

A queue always exposes its \emph{oldest} element at the front. A stack
needs to expose its \emph{newest} element at the top. If, every time we
push a new element, we immediately rotate the entire queue so that the
new element ends up at the front (instead of the back, where a plain
enqueue would place it), then the front of the queue always holds the most
recently pushed element -- exactly what a stack's \textit{top} needs.

\begin{ideabox}[Make Push Expensive, Keep Pop/Top Cheap]
$\textit{push}(x)$: enqueue $x$ normally (it lands at the back). Then
rotate the queue by dequeuing and re-enqueuing every \emph{other} element
(everything that was already in the queue before $x$ was added) -- this
moves $x$, one rotation step at a time, all the way around to the front.
$\textit{pop}()$ and $\textit{top}()$ then become trivial: simply dequeue
(or peek) the queue's front, which is always the most recently pushed
element by the invariant just established.
\end{ideabox}

\subsection*{Why rotating after every push maintains the LIFO invariant}

Each push performs exactly $(\textit{size before this push})$ rotation
steps, which is enough to move the newly enqueued element from the back
to the front while preserving the relative order of everything else
(each older element is dequeued and immediately re-enqueued once, so
their relative order among themselves never changes). After this
rotation, the queue's front-to-back order is exactly ``most recently
pushed first,'' i.e.\ LIFO order, by induction: if this invariant held
before the push, the rotation restores it after the push as well.

\subsection*{Algorithm}

\begin{enumerate}
  \item $\textit{push}(x)$: let $s$ be the queue's current size before
        this push. Enqueue $x$. Repeat $s$ times: dequeue an element and
        immediately re-enqueue it.
  \item $\textit{pop}()$: dequeue and return the front.
  \item $\textit{top}()$: return the front, without removing it.
\end{enumerate}

\subsection*{Dry run}

\dryrun{Push $1, 2, 3$ in order.}

\begin{itemize}
  \item Push $1$: queue $=[1]$ ($s=0$, no rotation needed).
  \item Push $2$: enqueue $2$: queue $=[1,2]$. Rotate $s=1$ time:
        dequeue $1$, enqueue $1$: queue $=[2,1]$.
  \item Push $3$: enqueue $3$: queue $=[2,1,3]$. Rotate $s=2$ times:
        dequeue $2$, enqueue $2$: queue $=[1,3,2]$; dequeue $1$, enqueue
        $1$: queue $=[3,2,1]$.
\end{itemize}

The queue's front-to-back order is $[3,2,1]$ -- exactly LIFO order, most
recently pushed ($3$) at the front.

\subsection*{C++ implementation}

\begin{lstlisting}
#include <bits/stdc++.h>
using namespace std;

class StackUsingQueue {
    queue<int> q;

public:
    void push(int x) {
        int s = q.size();
        q.push(x);
        for (int i = 0; i < s; i++) {
            q.push(q.front());
            q.pop();
        }
    }

    int pop() {
        int val = q.front();
        q.pop();
        return val;
    }

    int top() {
        return q.front();
    }
};

int main() {
    StackUsingQueue s;
    s.push(1); s.push(2); s.push(3);
    cout << "Top: " << s.top() << endl;
    cout << "Popped: " << s.pop() << endl;
    return 0;
}
\end{lstlisting}

\begin{complexitybox}
\textbf{Time Complexity.} $\textit{push}$ costs $O(n)$ (the rotation
touches every existing element). $\textit{pop}$ and $\textit{top}$ cost
$O(1)$.
\textbf{Space Complexity.} $O(n)$ for the underlying queue.
\end{complexitybox}

\begin{takeawaybox}
When forced to simulate one access order using a structure built for the
opposite order, you must pay an extra reordering cost \emph{somewhere} --
the only design choice is \emph{which} operation absorbs that cost. Here
we chose to make \textit{push} expensive and \textit{pop}/\textit{top}
cheap; Section~\ref{sec:sq-queue-using-stack} makes the opposite choice
for the reverse simulation, illustrating both options.
\end{takeawaybox}
